% Part: model-theory
% Chapter: models-of-arithmetic
% Section: standard-models

\documentclass[../../../include/open-logic-section]{subfiles}

\begin{document}

\olfileid[fr]{mod}{mar}{stm}
\olsection{Modèles standards de l'arithmétique}

Le !!{language} de l'arithmétique~$\Lang{L_A}$ est destiné
à parler des nombres, plus précisément des entiers naturels.
Le modèle standard~$\Struct{N}$ joue donc un rôle particulier :
c'est celui que nous voulons décrire. En logique, cependant,
les propriétés structurelles nous intéressent souvent davantage
que les objets qui composent le domaine. Deux !!{structure}s
isomorphes partagent ces propriétés. Nous appellerons donc
également « standard » toute !!{structure} isomorphe à~$\Struct{N}$.

\begin{defn}
Une !!{structure} pour $\Lang{L_A}$ est \emph{standard}
si elle est isomorphe à~$\Struct{N}$.
\end{defn}

\begin{prop}
\ollabel{prop:standard-domain}
Si une !!{structure}~$\Struct{M}$ est standard, son domaine
est l'ensemble des valeurs des numéraux standards, c'est-à-dire
\[
  \Domain{M}=\Setabs{\Value{\num{n}}{M}}{n\in\Nat}.
\]
\end{prop}

\begin{proof}
Il est clair que
$\Value{\num{n}}{M}\in\Domain{M}$ pour tout~$n$.
Réciproquement, soit $x\in\Domain{M}$. Puisque
$\Struct{M}$ est standard, il existe un isomorphisme
$g\colon\Nat\to\Domain{M}$. Il préserve la valeur des termes,
donc
$g(n)=g(\Value{\num{n}}{N})=\Value{\num{n}}{M}$.
Comme $g$ est !!{surjective}, $x=g(n)$ pour un certain
$n\in\Nat$, et par conséquent
$x=\Value{\num{n}}{M}$.
\end{proof}

\begin{explain}
Dans une !!{structure} standard~$\Struct{M}$, les numéraux
standards $\num{0}$, $\num{1}$, $\num{2}$, \dots,
c'est-à-dire les termes $\Obj{0}$, $\Obj{0}'$,
$\Obj{0}''$, etc., nomment tous les !!{element}s du
domaine. Cela ne signifie pas que ces !!{element}s
\emph{soient} des entiers naturels : les numéraux permettent
simplement de les désigner comme ils désignent les éléments
de~$\Domain{N}$.
\end{explain}

\begin{prob}
Montrer que la réciproque de la
\olref[mod][mar][stm]{prop:standard-domain} est fausse :
donner une !!{structure}~$\Struct{M}$ telle que
$\Domain{M}=
\Setabs{\Value{\num{n}}{M}}{n\in\Nat}$
mais qui ne soit pas isomorphe à~$\Struct{N}$.
\end{prob}

\begin{prop}
\ollabel{prop:thq-standard}
Si $\Sat{M}{\Th{Q}}$ et si
$\Domain{M}=
\Setabs{\Value{\num{n}}{M}}{n\in\Nat}$,
alors $\Struct{M}$ est standard.
\end{prop}

\begin{proof}
Définissons $g\colon\Nat\to\Domain{M}$ par
$g(n)=\Value{\num{n}}{M}$. L'hypothèse sur le
domaine assure que $g$ est !!{surjective}.
Elle est aussi !!{injective} : si $n\neq m$,
alors $\Th{Q}\Proves\eq/[\num{n}][\num{m}]$.
Puisque $\Sat{M}{\Th{Q}}$, il s'ensuit que
$\Sat{M}{\eq/[\num{n}][\num{m}]}$ et donc
$g(n)\neq g(m)$.

Vérifions que cette bijection est un isomorphisme.
\begin{enumerate}
\item Puisque $\Assign{\Obj{0}}{N}=0$ et
  $\num{0}$ est le terme $\Obj{0}$,
  \[
    g(\Assign{\Obj{0}}{N})=g(0)
    =\Value{\num{0}}{M}=\Assign{\Obj{0}}{M}.
  \]
\item Dans $\Struct{N}$, le symbole $\prime$ désigne
  le successeur. Or $\num{n+1}$ est le terme
  $\num{n}'$, donc
  \[
    g(\Assign{\prime}{N}(n))
    =g(n+1)=\Value{\num{n+1}}{M}
    =\Assign{\prime}{M}(\Value{\num{n}}{M})
    =\Assign{\prime}{M}(g(n)).
  \]
\item Dans $\Struct{N}$, $+$ désigne l'addition.
  Puisque
  $\Th{Q}\Proves\num{n+m}=(\num{n}+\num{m})$,
  on a
  \begin{align*}
    g(\Assign{+}{N}(n,m))
      &=g(n+m)=\Value{\num{n+m}}{M}\\
      &=\Value{\num{n}+\num{m}}{M}\\
      &=\Assign{+}{M}(
          \Value{\num{n}}{M},\Value{\num{m}}{M})\\
      &=\Assign{+}{M}(g(n),g(m)).
  \end{align*}
\item La vérification de
  $g(\Assign{\times}{N}(n,m))
  =\Assign{\times}{M}(g(n),g(m))$
  est laissée en exercice.
\item Enfin,
  $\tuple{n,m}\in\Assign{<}{N}$ si et seulement
  si $n<m$. Si $n<m$, alors
  $\Th{Q}\Proves\num{n}<\num{m}$ et,
  comme $\Struct{M}$ est un modèle de~$\Th{Q}$,
  $\tuple{g(n),g(m)}\in\Assign{<}{M}$.
  Si $n\not<m$, la théorie $\Th{Q}$ démontre
  $\lnot\num{n}<\num{m}$; par conséquent,
  $\tuple{g(n),g(m)}\notin\Assign{<}{M}$.
  L'appartenance à la relation $<$ est donc
  préservée dans les deux sens.
\end{enumerate}
\end{proof}

\begin{explain}
La fonction $g$ est la manière naturelle d'envoyer
$\Nat$ dans le domaine d'une !!{structure}~$\Struct{M}$
pour $\Lang{L_A}$ : chaque telle !!{structure} possède
des !!{element}s nommés par $\num{0}$, $\num{1}$,
$\num{2}$, etc. Si $\Struct{M}$ satisfait les
assertions fondamentales sur ces numéraux qui sont
vraies dans~$\Struct{N}$, et si $g$ est bijective,
elle devient un isomorphisme. Lorsque les numéraux
nomment tous les !!{element}s de~$\Domain{M}$,
cet isomorphisme est unique.
\end{explain}

\begin{prop}
\ollabel{prop:thq-unique-iso}
Si $\Struct{M}$ est standard, la fonction $g$ de
la preuve de la \olref{prop:thq-standard} est le seul
isomorphisme de $\Struct{N}$ sur~$\Struct{M}$.
\end{prop}

\begin{proof}
Soit $h\colon\Nat\to\Domain{M}$ un isomorphisme.
Montrons par récurrence sur~$n$ que $g(n)=h(n)$.
Au rang $0$, on a
$g(0)=\Assign{\Obj{0}}{M}$ par définition de~$g$;
comme $h$ est un isomorphisme,
$h(0)=h(\Assign{\Obj{0}}{N})
=\Assign{\Obj{0}}{M}$.

Supposons $g(n)=h(n)$. Alors
\begin{align*}
g(n+1)
  &=\Value{\num{n+1}}{M}
     &&\text{par définition de $g$}\\
  &=\Value{\num{n}'}{M}
     &&\text{car $\num{n+1}\ident\num{n}'$}\\
  &=\Assign{\prime}{M}(\Value{\num{n}}{M})
     &&\text{par définition de la valeur d'un terme}\\
  &=\Assign{\prime}{M}(g(n))
     &&\text{par définition de $g$}\\
  &=\Assign{\prime}{M}(h(n))
     &&\text{par hypothèse de récurrence}\\
  &=h(\Assign{\prime}{N}(n))
     &&\text{car $h$ est un isomorphisme}\\
  &=h(n+1).
\end{align*}
\end{proof}

\begin{explain}
Tout ensemble !!{denumerable}~$M$ est en bijection
avec~$\Nat$ et peut donc servir de domaine à un
modèle standard~$\Struct{M}$. Pour en obtenir un,
on choisit $z\in M$ comme valeur de $\Obj{0}$,
puis une fonction $s$ pour interpréter $\prime$.
L'orbite $z,s(z),s^2(z),\dots$ doit être
\emph{tout}~$M$, sans répétition; l'isomorphisme
$n\mapsto s^n(z)$ fixe alors les interprétations
de $+$, de $\times$ et de~$<$. En particulier,
$s$ est une bijection de $M$ sur
$M\setminus\{z\}$. Cette seule condition de
bijection ne suffit toutefois pas : une autre
orbite de $s$, disjointe de celle de~$z$,
donnerait une !!{structure} non standard.\footnote{Précision éditoriale : la source n'impose ici que les propriétés d'injectivité et d'image du successeur. Une orbite supplémentaire peut encore les satisfaire. L'orbite de $z$ doit donc épuiser le domaine pour que la structure construite soit standard ; l'image de $s$ est $M\setminus\{z\}$, et non $M$.}

La valeur $z$ ne peut appartenir à l'image de~$s$,
car $\lforall[x][\eq/[\Obj 0][x']]$ est vrai
dans~$\Struct{N}$. La fonction $s$ doit être
!!{injective}, comme le successeur dans~$\Struct{N}$;
elle doit aussi atteindre chaque élément de
$M\setminus\{z\}$, puisque
$\lforall[x][(\eq[x][\Obj 0]\lor
\lexists[y][\eq[y'][x]])]$ est vrai dans
le modèle standard.
\end{explain}

\end{document}
